\problemname{Jaga råttor}

\noindent
Förra året lyckades ingen jaga ut jätteråttorna ur Göteborg (se problemet \textit{KATT och råtta}).
På grund av detta har Olle inget annat val än att ta saken i egna händer och själv jaga in alla göteborgs
råttor i fällor.

Göteborgs kloaker består av $N$ brunnar, där $M$ stycken brunnar är direkt sammankopplade av rör. 
När Olle ska jaga en råtta sätter han en fälla i brunn $T$, han börjar han i brunn $S$ och råttan börjar
i brunn $R$. Målet är att få råttan att gå in i fällan i brunn $T$. Olle kan förflytta sig mellan brunnarna
med hjälp av rören. Om han går in i samma brunn som råttan, så kommer den fly till en intilliggande brunn.
Efter att ha jagat många råttor har han märkt att de är förutsägbara: om han går in i en brunn från ett
visst rör, så kommer råttan alltid att fly till en särskild brunn.

Olle har tyvärr märkt att vissa råttor inte går att fånga. Han vill därför ha din hjälp $Q$ gånger att avgöra om det går
att fånga råtten givet $S,T$ och $R$. Notera att alla frågor är oberoende varandra, och det finns alltid
exakt en rådda i kloakerna per fråga.




\section*{Indata}
Den första raden i indatan innehåller heltalen $N, M$ och $Q$ ($3 \leq N \leq 10^5$, $1 \le M,Q \le 10^5$),
antalet brunnar, antalet rör och antalet frågor.

De följande $M$ raderna beskriver vardera ett rör. Varje rör beskrivs av heltalen $u,v,a,b$ 
($1 \leq u,v \leq N$, $u \neq v$), där $u,v$ beskriver att det finns ett rör mellan brunn $u$ och $v$.
$a$ beskriver att om Olle färdas från $u$ till $v$ och en råtta befinner sig i brunn $v$, så kommer 
den fly till brunn $a$. $b$ beskriver motsvarande att om Olle befinner sig i brunn $v$ och färdas
till $u$, så flyr råtten till brunn $b$. Det är även garanterat att det finns som mest ett rör
mellan varje par av brunnar, röret $(v,a)$ finns och röret $(u,b)$ finns.

Därefter följer $Q$ rader som vardera innehåller heltalen $S,R,T$ ($1 \leq S,R,T \leq N$),
som beskriver att Olle undrar om han kan jaga råttan i brunn $R$ till brunn $T$ om han börjar på
brunn $S$. Det är garanterat att $T\neq S$, $T\leq R$ och $S\neq R$.


\section*{Utdata}
Svara på varje fråga i samma ordning som i indatan. Om Olle kan jaga in råttan i fällen, skriv ut
\texttt{1}, annars \textit{0}.

\section*{Poängsättning}
Din lösning kommer att testas på en mängd testfallsgrupper.
För att få poäng för en grupp så måste du klara alla testfall i gruppen.

%TODO
% föreställer mig att lägga till låggradigt + a och b är lika med grannen med minst index
% <1, M^2Q>
% <M^2, MQ>
% <N*maxdeg^2*log,Q*M*maxdeg^2>
% träd
% varje a och b pekar på giltiga grannen med lägst index
\noindent
\begin{tabular}{| l | l | p{12cm} |}
  \hline
  \textbf{Grupp} & \textbf{Poäng} & \textbf{Gränser} \\ \hline
  $1$    & $5$        & $N, M \leq 50, Q \leq 100$ \\ \hline
  $2$    & $11$       & $N, M \leq 500, Q \leq 1000$ \\ \hline
  $3$    & $26$       & $Q=1$, $deg(u) \leq 5$ för alla $1 \leq u \neq N$. \\ \hline
  $4$    & $25$       & $Q \leq 100$ \\ \hline
  $5$    & $7$        & $R$ och $T$ är samma för alla frågor \\ \hline
  $6$    & $11$        & $T$ är samma för alla frågor \\ \hline
  $7$    & $15$       & Inga ytterligare begränsningar. \\ \hline
\end{tabular}

\section*{Förklaring av exempelfall}
